% ============================================================================
%  5. DESAF\'IOS DE IMPLEMENTACI\'ON Y SOLUCIONES APLICADAS
%  ----------------------------------------------------------------------------
%  Los desafíos descritos corresponden a dificultades reales
%  documentadas durante la ejecución de las actividades a4, a5 y a7.
% ============================================================================

\section{Desaf\'ios de implementaci\'on y soluciones aplicadas}

La puesta en marcha de la infraestructura plante\'o dificultades reales
que no se anticipate durante la planificaci\'on inicial. Este cap\'itulo
recoge los cinco desaf\'ios m\'as significativos, describe su sintomatizaci\'on
y expone la soluci\'n adoptada por el grupo, incluyendo las verificaciones
que permitieron confirmar la resoluci\'on.

\subsection{Conflicto entre dos servidores \textsc{dhcp} activos}

El primer obstacle se manifest\'o al conectar el primer cliente a la red:
obten\'ia una direcci\'n de \texttt{192.168.0.1}, correspondiente al router,
en lugar de recibirla del servidor Kea del proyecto. La causa era que el
router ten\'ia su propio servicio \textsc{dhcp} activo, de modo que tanto
\'este como Kea respond\'ian a los paquetes \textsc{discover} con ofertas
distintas. El cliente aceptaba la primera oferta recibida, lo que
provocaba una configuraci\'n de red inconsistente y, en consecuencia,
direcciones duplicadas entre equipos \cite{rfc2131}.

Para resolverlo se accedi\'o a la interfaz de administraci\'on del router y
se deshabilit\'o el servicio \textsc{dhcp} en la \textsc{lan}, manteniendo
\'unicamente la configuraci\'on de la puerta de enlace y las reglas de
red. La figura \ref{fig:desafio-dhcp} documenta el estado final de dicho
par\'ametro.

\capturaAPA{fig:desafio-dhcp}{Desactivaci\'on del servicio \textsc{dhcp} del router para eliminar el conflicto con Kea}{captura-router-disable.png}{Captura de pantalla de la configuraci\'on \textsc{lan} del router TP-Link donde la opci\'on \textsc{dhcp server} figura en estado deshabilitado, resolviendo la Respuesta duplicada de los dos servidores \textsc{dhcp}.}

\subsection{Verificaci\'on previa de la conectividad del cliente}

Tras la desactivaci\'on del \textsc{dhcp} del router se comprob\'o que
ning\'un servicio de la red hubiera quedado inaccesible como efecto
secundario. Esta verificaci\'on result\'o relevante porque el router
concentra la puerta de enlace y la salida a Internet: si alg\'un equipo
hab\'ia quedado sin ruta por defecto, el cambio habr\'ia afectado a todos
los equipos que dependieran de ella. La figura \ref{fig:desafio-ping}
muestra la verificaci\'on realizada.

\capturaAPA{fig:desafio-ping}{Verificaci\'on de la conectividad del cliente hacia la puerta de enlace tras el cambio}{captura-ping-router-gateway.png}{Captura de pantalla de la terminal del cliente con las respuestas de \texttt{ping} hacia la puerta de enlace \texttt{192.168.0.1}, confirmando que el equipo mantiene su conectividad tras la intervenci\'on.}

\subsection{Resoluci\'on err\'onea del \textit{hostname} a \textit{loopback}}

Al instalar el controlador de dominio \texttt{dc1} se observ\'o que las
operaciones de uni\'on al dominio fallaban con errores de acceso al
servidor Kerberos. El origen estaba en la resoluci\'on del nombre del host:
si \texttt{dc1} se resuelve a \texttt{127.0.0.1}, el servidor se comunica
consigo mismo a trav\'es del \textit{loopback} y los clientes de la red
nunca alcanzan al controlador de dominio \cite{sambateam2024}.

La soluci\'n consisti\'o en garantizar que el nombre \texttt{dc1} resuelva a
la direcci\'n \texttt{192.168.0.2/24} de la interfaz cableada \texttt{eno1}
y registrar esa correspondencia de forma explícita en \texttt{/etc/hosts},
de modo que la resoluci\'o local tenga prioridad sobre cualquier otra
entrada.

\capturaAPA{fig:desafio-loopback}{Correspondencia entre el nombre del host y su direcci\'n de red en \texttt{/etc/hosts}}{captura-ip-fija.png}{Captura de pantalla de la salida de \texttt{ip -4 a show eno1} junto con la entrada de \texttt{/etc/hosts} que asocia \texttt{dc1.sudoers.lan} a la direcci\'n de la red local y no al \textit{loopback}.}

\subsection{Desfase de reloj y autenticaci\'on Kerberos}

Las lecturas de tickets de Kerberos fallaban de forma intermitente en los
clientes, con el error de versi\'on de ticket incorrecta. La causa no era la
configuraci\'on del dominio sino el reloj del sistema: Kerberos exige que
el desfase entre el cliente y el servidor sea inferior a cinco minutos
para aceptar el ticket \cite{rfc4120}.

La soluci\'on fue instalar y habilitar el servicio \texttt{chrony} en el
host del controlador de dominio, de modo que la hora del servidor se
sincronizara de forma continua con los servidores de referencia. Tras
activar la sincronizaci\'on, la obtenci\'on de tickets se complet\'o de
forma satisfactoria.

\capturaAPA{fig:desafio-kerberos}{Servicios de sincronizaci\'on horaria habilitados en el host del dominio}{captura-servicios-host.png}{Captura de pantalla de la salida de \texttt{systemctl is-enabled} mostrando \texttt{chrony} junto con \texttt{samba-ad-dc}, \texttt{nftables} y \texttt{docker} habilitados en el arranque del servidor.}

\subsection{Concesi\'on \textsc{dhcp} heredada del router}

Tras resolver el conflicto de servidores \textsc{dhcp}, los equipos que ya
ten\'ian una direcci\'n asignada por el router la conservaban, ya que la
concesi\'on anterior segu\'ia siendo v\'alida. Estos clientes no
actualizaban su configuraci\'on de red y, en consecuencia, persist\'an sin
poder resolver los nombres del dominio \cite{droms1997}.

La soluci\'on consisti\'o en forzar la renovaci\'on de la concesion desde
el cliente, de modo que \'este la solicite al servidor Kea y reciba una
direcci\'n dentro del \'ambito din\'amico junto con la puerta de enlace y el
servidor \textsc{dns} correctos. Esta renovaci\'n tambi\'n sirvi\'o como
prueba de que el servidor \textsc{dhcp} del proyecto asum\'ia \'unicamente
la totalidad de las concesiones de la red.

\capturaAPA{fig:desafio-dhclient}{Renovaci\'on forzada de la concesion \textsc{dhcp} en el cliente}{captura-dhclient.png}{Captura de pantalla de la terminal del cliente mostrando la renovaci\'on de la direcci\'n y la asignaci\'n de una direcci\'n dentro del \'ambito din\'amico definido por Kea.}

\subsection{Servicios que requieren red \textit{host} en contenedores}

Por \'ultimo, el despliegue del servicio \textsc{dhcp} dentro de Docker
present\'o un inconveniente de conectividad: el protocolo utiliza la
direcci\'on de difusi\'on por \textsc{broadcast} en el puerto
\texttt{67/udp}, y los contenedores con red \textit{bridge} no reenv\'ian
ese tipo de tr\'afico. Como resultado, los clientes no recib\'ian oferta
alguna.

La soluci\'on fue configurar el contenedor con la opci\'n
\texttt{network\_mode: host}, que descarta el aislamiento de red y permite
al servicio escuchar directamente en la interfaz del host. Esta decisi\'on
se tom\'o de forma consciente, asumiendo que el servicio queda expuesto a
la \textsc{lan} y que por ello debe protegerse mediante las reglas de
firewall del host \cite{stewart2017}.

\subsection{Advertencia de seguridad del navegador por autoridad de certificaci\'on no reconocida}

Al publicar el portal por \textsc{https} con el certificado
\textit{wildcard} \texttt{*.sudoers.lan}, los navegadores de los clientes
mostraban una advertencia de que la conexi\'on no era privada. El servidor
respond\'a correctamente y el cifrado es real, pero el certificado lo firma
una autoridad de certificaci\'on interna generada por el propio grupo y, por
tanto, no figura en la lista de autoridades en las que el navegador conf\'ia
por defecto. Para el usuario esto se manifiesta como un riesgo inexistente
que induce a ignorar la advertencia.

La soluci\'on consisti\'o en distribuir el certificado de la autoridad
(\texttt{ca.crt}) e instalarlo como entidad de confianza en el sistema de
cada estaci\'on. Una vez instalada, el candado aparece en la barra de
direcciones sin advertencias y el cifrado queda validado de extremo a
extremo. La figura \ref{fig:desafio-tls} documenta el acceso al portal una
vez resuelta la cadena de confianza.

\capturaAPA{fig:desafio-tls}{Acceso al portal tras instalar la autoridad de certificaci\'on interna en el cliente}{captura-cliente-web.png}{Captura de pantalla del navegador de un cliente mostrando el portal de la intranet con el candado de confianza de la autoridad interna, sin la advertencia de certificado no reconocido.}
